\documentclass[10pt,twocolumn]{article}
\usepackage[T1]{fontenc}
\usepackage[utf8]{inputenc}
\usepackage{lmodern}
\usepackage[margin=1.8cm,top=2.4cm,bottom=2cm,columnsep=0.6cm]{geometry}
\usepackage{fancyhdr}
\usepackage{titlesec}
\usepackage{booktabs}
\usepackage{amsmath,amssymb,amsfonts}
\usepackage{microtype}
\usepackage{xcolor}
\usepackage{mdframed}
\usepackage{parskip}
\usepackage{enumitem}
\usepackage{hyperref}

\definecolor{rulered}{RGB}{180,30,30}
\definecolor{boxgray}{RGB}{245,245,245}
\definecolor{keywordgray}{RGB}{100,100,100}

\pagestyle{fancy}
\fancyhf{}
\fancyhead[L]{\textsc{\small Claude Mag}}
\fancyhead[C]{\small\textit{Machine Learning Research Digest}}
\fancyhead[R]{\small 2026-09-28}
\fancyfoot[C]{\small\thepage}
\renewcommand{\headrulewidth}{0.8pt}

\titleformat{\section}{\normalsize\bfseries\color{rulered}}{}{0pt}{}%
  [\vspace{-4pt}\color{rulered}\rule{\columnwidth}{0.4pt}]
\titlespacing{\section}{0pt}{8pt}{4pt}

\hypersetup{colorlinks=true,urlcolor=rulered,linkcolor=black}

\begin{document}
\setlength{\parindent}{0pt}
\setlength{\parskip}{4pt}

{\small\color{keywordgray}\textbf{Keywords:}
  robot manipulation \textbullet{}
  LLM planning \textbullet{}
  adaptive control \textbullet{}
  physical state}\par\vspace{4pt}

{\LARGE\bfseries\raggedright
  Body-Grounded Replanning}\par\vspace{4pt}

{\large\itshape\raggedright
  Physically Adaptive Robot Manipulation via
  LLM-Based Internal State Replanning}\par\vspace{6pt}

{\small\textbf{By} Namiko Saito \& Hiroshi Kera
  (Microsoft Research Asia Tokyo; Waseda University;
  Chiba University; National Institute of Informatics)}\par\vspace{2pt}

{\small\color{keywordgray}arXiv:2609.30024 \textbullet{} September 2026}
\par\vspace{2pt}\noindent{\color{rulered}\rule{\columnwidth}{0.6pt}}\vspace{6pt}

Body-grounded high-level replanning uses internal physical signals---joint
torques and mobility constraints---to trigger LLM-based strategy switching
during manipulation, maintaining task success while substantially reducing
physical effort under load and asymmetric mobility conditions.

\begin{mdframed}[backgroundcolor=boxgray,linecolor=rulered,linewidth=1pt,
    innerleftmargin=6pt,innerrightmargin=6pt,
    innertopmargin=6pt,innerbottommargin=6pt]
\textbf{\small AT A GLANCE}\par\vspace{4pt}
\begin{description}[leftmargin=0pt,labelsep=4pt,itemsep=2pt,parsep=0pt,
    font=\small\bfseries]
  \item[Problem:] A manipulation strategy may remain geometrically feasible
    while becoming physically unsuitable due to increased joint load or
    limited mobility; current planners have no mechanism to detect or
    respond to this.
  \item[Why it matters:] Robots deployed in real environments accumulate
    physical state changes that degrade performance and increase wear; without
    adaptation, task success and hardware longevity both suffer.
  \item[Method:] Body-state events trigger an LLM that interprets joint-level
    state and execution history to select a physically appropriate alternative
    strategy from a library, without changing the task objective or low-level
    controller.
  \item[Result:] High task success maintained; physical effort reduced; more
    efficient strategy adaptation than without body state awareness.
  \item[Evidence:] Reaching task under controlled load and asymmetric mobility
    constraints, simulation and real robot.
\end{description}
\end{mdframed}

\section{The Big Picture}

Robotic manipulation planners reason carefully about the external world:
obstacle geometry, object poses, grasp feasibility, and collision avoidance.
But they typically treat the robot's own body as a fixed, ideal actuator
whose physical state does not change during execution.

In practice, this assumption fails. A joint under sustained load accumulates
fatigue. A component subject to hardware wear constrains the workspace
asymmetrically. A strategy that was geometrically and dynamically feasible
at planning time may become physically unsuitable minutes later. The standard
response is to run the low-level controller harder---increasing torque
outputs to compensate---but this accelerates wear, increases energy
consumption, and can ultimately cause failure.

Body-grounded replanning introduces a new level in the robot's deliberative
hierarchy: the ability to recognize when internal physical state has made
the current high-level strategy inappropriate, and to select a physically
better alternative without changing the task objective or the low-level
controller.

\section{Why Existing Approaches Fall Short}

\textbf{Standard task-and-motion planning (TAMP)} optimizes over geometry
and kinematics. Joint torque, fatigue, and mobility constraints are not
modeled at the planning level; they are delegated to the controller.

\textbf{Adaptive controllers} adjust joint-level torques and impedances
in response to physical state but do not change high-level strategy (e.g.,
which direction to approach, which grasp type to use).

\textbf{LLM-based planners} condition on task descriptions and environmental
observations---scene graphs, object attributes, occupancy maps---but have no
access to joint-level physical signals.

\textbf{Fault-tolerant planning} handles discrete, identifiable hardware
failures (a joint locked at zero torque) but not continuous, graded physical
state degradation.

\section{Core Mechanism}

\textbf{Body-State Event Detection.} Physical state $s_t = (\tau_t, q_t,
v_t)$ is monitored continuously. An event triggers when:
\begin{itemize}[itemsep=2pt,parsep=0pt]
  \item Joint load: $\max_j \tau_{j,t} > \tau_\text{thresh}$
  \item Mobility constraint: reachability $R(q_t) < R_\text{thresh}$
\end{itemize}

\textbf{LLM Strategy Selector.} On event, the LLM receives a prompt encoding:
\begin{enumerate}[itemsep=2pt,parsep=0pt]
  \item \textbf{Joint-level state:} current torques and a natural-language
    summary of load distribution and mobility limits.
  \item \textbf{Execution statistics:} success rate and mean completion time
    of recent strategy attempts.
  \item \textbf{Execution history:} which strategies were tried and with
    what outcomes.
\end{enumerate}

The LLM selects a new strategy from the strategy library (e.g., ``approach
from the side,'' ``use extended-arm reach,'' ``shift load to stronger joints'').

\textbf{Decoupled Architecture.} High-level replanning and low-level
execution are fully decoupled: replanning changes the strategy but not the
task objective or the controller's internal parameters.

\section{Technical Formulation}

Let $\sigma \in \mathcal{S}$ be the current strategy from library $\mathcal{S}$.
Event condition at time $t$:
\[
e_t = \mathbf{1}\!\left[
  \max_j \tau_{j,t} > \tau_\text{thresh}
  \;\vee\;
  R(q_t) < R_\text{thresh}
\right]
\]

On $e_t = 1$, the LLM is queried:
\[
\sigma_{t+1} = \mathrm{LLM}\!\left(
  \mathcal{P}(s_{t-T:t},\; h_{0:t},\; \sigma_\text{hist})
\right)
\]

where $\mathcal{P}$ encodes physical state over window $T$, execution
history $h$, and prior strategies $\sigma_\text{hist}$.

The optimization objective:
\[
\max_{\sigma_{0:\infty}} \Pr[\text{task success}]
\;\;\text{s.t.}\;\;
\sum_t \|\tau_t\|_2 \leq E_\text{budget}
\]

\section{Learning or Inference Procedure}

No model training is required. At execution time:
\begin{enumerate}[itemsep=2pt,parsep=0pt]
  \item Execute strategy $\sigma$ with the low-level controller.
  \item Monitor $s_t$ at each control step.
  \item On $e_t = 1$: construct prompt $\mathcal{P}$, query the LLM.
  \item Update $\sigma \leftarrow \sigma_{t+1}$ and resume.
  \item Repeat until success or timeout.
\end{enumerate}

The LLM runs at inference time only, with no task-specific fine-tuning.

\section{What the Guarantee Says}

The framework does not require training a new model. Its correctness relies
on: (i) the strategy library covering alternative strategies that are
physically better in the relevant conditions, and (ii) the LLM's ability to
interpret physical state descriptions and select appropriate alternatives.
Evaluation shows that both conditions hold for the tested configurations.

\section{Experimental Findings}

\begin{itemize}[itemsep=2pt,parsep=0pt]
  \item \textbf{Task:} Reaching under two constraint types---controlled
    joint load (added deadweight) and asymmetric mobility (reduced joint range)
  \item \textbf{Platforms:} Simulated and real robot arm
  \item Body-grounded replanning maintains high task success rates where
    fixed-strategy baselines fail under increasing load
  \item Cumulative joint torque is substantially reduced by selecting
    strategies that distribute load across available joints
  \item Strategy adaptation is more efficient (fewer failed attempts before
    success) than random strategy switching when mobility is asymmetric
  \item Performance is robust across multiple off-the-shelf LLMs; stronger
    models provide modestly better strategy selection
\end{itemize}

\section{Ablations and Interpretation}

\textbf{Event threshold:} Too low $\tau_\text{thresh}$ causes excessive
replanning; too high causes physical overload. Performance is robust within
a reasonable range.

\textbf{History window:} More execution history improves strategy selection
for mobility constraints; load constraints show less sensitivity.

\textbf{Strategy library size:} A larger library with more alternatives
improves success under diverse constraints; the LLM selects appropriately
among them.

\textbf{Architecture decoupling:} The clean separation between replanning and
execution means the framework can be added to any existing manipulation
system without modifying the low-level controller.

\section*{Reference}

Namiko Saito, Hiroshi Kera.
\textbf{Body-Grounded Replanning for Physically Adaptive Manipulation}.
arXiv:2609.30024, September 2026.\\
\url{https://arxiv.org/abs/2609.30024}

\end{document}
